% OpenStax Calculus Volume 1 -- Section 4.4 Exercises: The Mean Value Theorem
% Exercises reproduced from OpenStax, licensed under CC BY 4.0.
% Source: https://openstax.org/books/calculus-volume-1/pages/4-4-the-mean-value-theorem
\documentclass{exercises}
\title{OpenStax Calculus Volume 1}
\subtitle{Section 4.4 Exercises: The Mean Value Theorem}
\sourceurl{https://openstax.org/books/calculus-volume-1/pages/4-4-the-mean-value-theorem}
\author{}
\date{}

\begin{document}
\maketitle

\begin{enumerate}[start=1]
\item Why do you need continuity to apply the Mean Value Theorem? Construct a counterexample.
\item Why do you need differentiability to apply the Mean Value Theorem? Find a counterexample.
\item When are Rolle's theorem and the Mean Value Theorem equivalent?
\item If you have a function with a discontinuity, is it still possible to have $f'(c)(b-a)=f(b)-f(a)$? Draw such an example or prove why not.
\end{enumerate}

For the following exercises, determine over what intervals (if any) the Mean Value Theorem applies. Justify your answer.

\begin{enumerate}[resume]
\item $y=\sin(\pi x)$
\item $y=\frac{1}{x^{3}}$
\item $y=\sqrt{4-x^{2}}$
\item $y=\sqrt{x^{2}-4}$
\item $y=\ln(3x-5)$
\end{enumerate}

For the following exercises, graph the functions on a calculator and draw the secant line that connects the endpoints. Estimate the number of points $c$ such that $f'(c)(b-a)=f(b)-f(a)$.

\begin{enumerate}[resume]
\item \techrequired $y=3x^{3}+2x+1$ over $[-1,1]$
\item \techrequired $y=\tan(\frac{\pi}{4}x)$ over $[-\frac{3}{2},\frac{3}{2}]$
\item \techrequired $y=x^{2}\cos(\pi x)$ over $[-2,2]$
\item \techrequired $y=x^{6}-\frac{3}{4}x^{5}-\frac{9}{8}x^{4}+\frac{15}{16}x^{3}+\frac{3}{32}x^{2}+\frac{3}{16}x+\frac{1}{32}$ over $[-1,1]$
\end{enumerate}

For the following exercises, find all points $0<c<2$ predicted by the Mean Value Theorem satisfying $f(2)-f(0)=f'(c)(2-0)$.

\begin{enumerate}[resume]
\item $f(x)=x^{3}$
\item $f(x)=\sin(\pi x)$
\item $f(x)=\cos(2\pi x)$
\item $f(x)=1+x+x^{2}$
\item $f(x)=(x-1)^{10}$
\item $f(x)=(x-1)^{9}$
\end{enumerate}

For the following exercises, show there is no $c$ such that $f(1)-f(-1)=f'(c)(2)$. Explain why the Mean Value Theorem does not apply over the interval $[-1,1]$.

\begin{enumerate}[resume]
\item $f(x)=|x-\frac{1}{2}|$
\item $f(x)=\frac{1}{x^{2}}$
\item $f(x)=\sqrt{|x|}$
\item $f(x)=\lfloor x\rfloor$ (Hint: This is called the floor function and it is defined so that $f(x)$ is the largest integer less than or equal to $x$.)
\end{enumerate}

For the following exercises, determine whether the Mean Value Theorem applies for the functions over the given interval $[a,b]$. Justify your answer.

\begin{enumerate}[resume]
\item $y=e^{x}$ over $[0,1]$
\item $y=\ln(2x+3)$ over $[-\frac{3}{2},0]$
\item $f(x)=\tan(2\pi x)$ over $[0,2]$
\item $y=\sqrt{9-x^{2}}$ over $[-3,3]$
\item $y=\frac{1}{|x+1|}$ over $[0,3]$
\item $y=x^{3}+2x+1$ over $[0,6]$
\item $y=\frac{x^{2}+3x+2}{x}$ over $[-1,1]$
\item $y=\frac{x}{\sin(\pi x)+1}$ over $[0,1]$
\item $y=\ln(x+1)$ over $[0,e-1]$
\item $y=x\sin(\pi x)$ over $[0,2]$
\item $y=5+|x|$ over $[-1,1]$
\end{enumerate}

For the following exercises, consider the roots of the equation.

\begin{enumerate}[resume]
\item Show that the equation $y=x^{3}+4x+16$ has exactly one real root. What is it?
\item Find the conditions for exactly one root (double root) for the equation $y=x^{2}+bx+c$
\item Find the conditions for $y=e^{x}-b$ to have one root. Is it possible to have more than one root?
\end{enumerate}

For the following exercises, use a calculator to graph the function over the interval $[a,b]$ and graph the secant line from $a$ to $b$. Use the calculator to estimate all values of $c$ as guaranteed by the Mean Value Theorem. Then, find the exact value of $c$, if possible, or write the final equation and use a calculator to estimate to four digits.

\begin{enumerate}[resume]
\item \techrequired $y=\tan(\pi x)$ over $[-\frac{1}{4},\frac{1}{4}]$
\item \techrequired $y=\frac{1}{\sqrt{x+1}}$ over $[0,3]$
\item \techrequired $y=|x^{2}+2x-4|$ over $[-4,0]$
\item \techrequired $y=x+\frac{1}{x}$ over $[\frac{1}{2},4]$
\item \techrequired $y=\sqrt{x+1}+\frac{1}{x^{2}}$ over $[3,8]$
\item At 10:17 a.m., you pass a police car at 55 mph that is stopped on the freeway. You pass a second police car at 55 mph at 10:53 a.m., which is located 39 mi from the first police car. If the speed limit is 60 mph, can the police cite you for speeding?
\item Two cars drive from one stoplight to the next, leaving at the same time and arriving at the same time. Is there ever a time when they are going the same speed? Prove or disprove.
\item Show that $y=\sec^{2}x$ and $y=\tan^{2}x$ have the same derivative. What can you say about $y=\sec^{2}x-\tan^{2}x$?
\item Show that $y=\csc^{2}x$ and $y=\cot^{2}x$ have the same derivative. What can you say about $y=\csc^{2}x-\cot^{2}x$?
\end{enumerate}

\end{document}
